\section{问题分析}

本题的核心是解释封锁冲击如何经过供需缺口、政策缓冲、运输替代、库存消耗和风险预期传导到布伦特原油价格。题面给定的短期低弹性会放大无缓冲反事实价格，但附件真实价格只在约 110--120 美元/桶附近形成平台，说明市场中存在多重缓冲和再定价过程。因此，本文采用“反事实基准--短期递推--中长期情景--敏感性检验”的建模路线。

\begin{enumerate}
  \item 传统供需基准用于估计无缓冲条件下的理论压力上界，为后续模型提供反事实参照。
  \item 短期冲击模型在日度递推中纳入供应中断、SPR、商业库存、绕道运输、需求收缩、恐慌衰减和预期修复，用于解释附件冲突窗口内的真实价格路径。
  \item 中长期油价调节模型在 60--180 天范围内设置乐观、中性和悲观三类情景，并用蒙特卡洛情景树和状态转移模型刻画尾部风险。
  \item 敏感性分析和外部风险校验用于识别关键驱动因素，避免把单一路径误读为确定性预测。
\end{enumerate}

\begin{figure}[htbp]
\centering
\begin{tikzpicture}[
  node distance=0.82cm and 0.72cm,
  stage/.style={draw=OilBorder, fill=PaperNoteBack, rounded corners=1pt, align=center, minimum height=0.78cm, text width=2.85cm, font=\small\bfseries},
  note/.style={draw=OilBorder, fill=white, rounded corners=1pt, align=center, minimum height=0.62cm, text width=2.85cm, font=\scriptsize},
  arrow/.style={-{Latex[length=2.2mm]}, line width=0.45pt, draw=OilGray}
]
\node[stage] (data) {附件价格数据};
\node[stage, right=of data] (base) {传统供需基准};
\node[stage, right=of base] (short) {短期冲击模型};
\node[stage, right=of short] (long) {中长期油价调节模型};
\node[note, below=of base] (counter) {低弹性反事实压力};
\node[note, below=of short] (check) {基准对比\\消融与稳健性};
\node[note, below=of long] (risk) {三情景\\蒙特卡洛与状态转移};
\draw[arrow] (data) -- (base);
\draw[arrow] (base) -- (short);
\draw[arrow] (short) -- (long);
\draw[arrow] (base) -- (counter);
\draw[arrow] (short) -- (check);
\draw[arrow] (long) -- (risk);
\draw[arrow] (counter.east) -- (check.west);
\draw[arrow] (check.east) -- (risk.west);
\end{tikzpicture}
\caption{本文总体技术路线}
\label{fig:paper-roadmap}
\end{figure}

这一技术路线的核心原则是：真实价格只来自附件数据；长期部分采用情景预测和区间表达；模型检验覆盖朴素基准、滞后复制、参数过拟合和结构稳定性等关键边界。

最终论文使用赛题给定的两个模型名称组织主线：任务 1 为短期冲击模型，任务 2 为中长期油价调节模型。传统供需模型用于提供反事实基准；三情景、蒙特卡洛和状态转移情景树是中长期油价调节模型内部的条件外推方法；消融实验、敏感性分析和蒙特卡洛情景树用于证明因素是否值得保留，并刻画多因素同时恶化时的尾部风险。因素纳入原则见\cref{tab:factor-inclusion}。

\begin{table}[!htbp]
\centering
\caption{因素纳入原则}
\label{tab:factor-inclusion}
\begin{tabularx}{\linewidth}{lX X}
\toprule
层级 & 因素 & 处理方式 \\
\midrule
短期冲击模型物理层 & 供应中断、短期弹性、SPR、商业库存、绕道运输、需求收缩、恐慌衰减 & 进入短期冲击模型，参与 0--60 天路径拟合与机制解释 \\
短期冲击模型价格形成层 & 地缘风险溢价、冲击不确定性、持续封锁制度风险、缓冲确认折价、预期修复折价 & 进入短期冲击模型，但必须通过基准对比、消融实验和稳健性检验证明必要性 \\
中长期调节层 & OPEC+ 剩余产能、非 OPEC 供给响应、冲突升级概率 & 不强行用于短期拟合，进入中长期油价调节模型和尾部风险解释 \\
改进方向 & 油轮保险费、美元指数、期货期限结构、投机资金和套保需求 & 当前缺少附件内真实外生数据，不单独参数化，作为后续扩展方向 \\
\bottomrule
\end{tabularx}
\end{table}

\keypoint{本文围绕三个问题展开：传统供需模型为什么给出高理论压力；哪些缓冲机制能把价格压回现实平台；封锁延长至 60--180 天时，哪些因素决定价格中枢和尾部跳涨风险。}

\paragraph{文献依据与建模约束}

相关研究表明，油价冲击不能只由当期供应缺口解释，还需要考虑库存、预期、地缘风险和宏观需求环境。Kilian（2009）、Hamilton（2009）和 Kilian 与 Murphy（2014）的研究说明，库存与预防性需求会改变冲击传导路径。因此，本文将传统供需模型定位为低弹性反事实上界，而不是最终预测模型。

在预测评价方面，Alquist 与 Kilian（2010）以及 Alquist、Kilian 与 Vigfusson（2013）指出，原油价格预测需要与无变化预测、随机游走等基准比较，并明确预测期限和不确定性边界。因此，本文在短期部分设置上一日价格朴素基准、漂移随机游走和滚动 ARIMA；在长期部分同时给出三情景中心路径、蒙特卡洛扇形区间和状态转移尾部风险。

在政策缓冲方面，Newell 与 Prest（2017）、Kilian 与 Zhou（2019）和 Stevens 与 Zhang（2021）均表明，SPR 释放能够影响油价，但其作用更接近削峰和稳定预期，而不是完全抵消供应中断。因此，本文把 SPR、商业库存和绕道运输都作为部分缓冲变量：它们改变剩余缺口和价格峰值，却不直接决定长期均衡价格。

在风险变量方面，Caldara 与 Iacoviello（2022）提供了 GPR 指数的理论和数据基础，Mei 等（2020）、Liu 等（2019）说明地缘风险更适合解释油价波动和尾部风险。本文将 GPR 用于外部风险校验和滞后检验，将 Cboe/FRED 的 OVX 原油 ETF 隐含波动率用于长期风险权重和不确定性强度约束，而不把二者作为短期同步拟合因子。

\begin{table}[!htbp]
\centering
\caption{核心文献对本文模型设计的约束}
\begin{tabularx}{\linewidth}{lX X}
\toprule
文献主线 & 对本文的约束 & 落地位置 \\
\midrule
油价冲击分解 & 不能只用供应缺口解释油价，需纳入库存、预期和风险重估 & 传统基准与短期冲击模型 \\
随机游走强基准 & RMSE 必须和 no-change、Random Walk、ARIMA 等基准比较 & 短期模型基准对比与 DM 检验 \\
SPR 与库存政策缓冲 & SPR 是削峰和争取时间，不是完全填补缺口的万能变量 & 物理供需层、消融实验和长期敏感性分析 \\
GPR 与 OVX 波动风险 & GPR 适合做风险校验，OVX 适合做市场隐含波动率约束，二者均不宜同步解释同日或同月油价方向 & 滞后 GPR 检验、OVX 滞后波动检验和长期尾部风险解释 \\
长期预测不确定性 & 中长期结果应报告情景、区间和尾部概率，不夸大单点预测 & 中长期油价调节模型、蒙特卡洛和状态转移情景树 \\
\bottomrule
\end{tabularx}
\end{table}
